\documentclass[11pt]{article}
\usepackage{amsmath,amssymb,amsthm,hyperref}
\usepackage[margin=1in]{geometry}
\newtheorem{theorem}{Theorem}
\title{A constant-map counterexample to a proposed chaos criterion}
\author{Principia Math In Progress}
\date{29 September 2026}
\begin{document}\maketitle
\begin{abstract}
The inequality $|\alpha+1|>|\beta|$ does not imply the existence of
chaotic solutions of $z_{n+1}=(\alpha+\alpha z_n+\beta z_{n-1})/(1+z_n)$.
For $\beta=0$ and $\alpha\ne-1$, every admissible solution is exactly
constant after one scalar step. In particular $\alpha=i$, $\beta=0$
disproves MathDB 331783 and the source's Conjecture 5.1.
\end{abstract}

\section*{Task 1: full counterexample and source scope}
Hassan's manuscript \cite{Hassan}, page 1, permits arbitrary complex
parameters and initial values. Conjecture 5.1 on page 13 asserts that
chaotic solutions exist if and only if $|\alpha+1|>|\beta|$.
The conjecture does not require $\beta\ne0$. The source identifies
chaos numerically through positive Lyapunov exponents on page 12.
The example below has a constant second iterate as a state-space map,
so it has no chaotic orbit under that criterion or under the usual
orbit-separation criteria.

\begin{theorem}
For every $\alpha\in\mathbb C\setminus\{-1\}$ and $\beta=0$, each
initial pair $(z_{-1},z_0)\in\mathbb C^2$ with $z_0\ne-1$ gives an
everywhere-defined forward solution satisfying
\[
 z_n=\alpha\qquad\text{for all }n\ge1.
\]
The strict inequality $|\alpha+1|>|\beta|$ holds, but there are no
chaotic solutions. In particular $(\alpha,\beta)=(i,0)$ is a
counterexample to the conjecture.
\end{theorem}
\begin{proof}
At every point where the original denominator is nonzero,
\[
 \frac{\alpha+\alpha z_n+0z_{n-1}}{1+z_n}
 =\frac{\alpha(1+z_n)}{1+z_n}=\alpha.
\]
Thus $z_1=\alpha$. Since $\alpha\ne-1$, the next denominator is
$1+\alpha\ne0$, and induction gives the asserted solution for all
$n\ge1$. No assignment at the excluded pole $z_n=-1$ is used.

For completeness, put $D=\{(x,y)\in\mathbb C^2:y\ne-1\}$.
The state-space map on this forward-invariant domain is
\[
 F(x,y)=(y,\alpha),\qquad F^2(x,y)=(\alpha,\alpha).
\]
Every pair of forward orbits therefore agrees exactly after two
iterations. In particular there is no persistent separation of nearby
orbits and no Li--Yorke scrambled pair: the distance between any two
iterated states is zero for every $n\ge2$.
The complex Jacobian is
\[
 DF=\begin{pmatrix}0&1\\0&0\end{pmatrix},\qquad
 D(F^n)=0\quad(n\ge2).
\]
Viewed as a real four-dimensional map, the same nilpotence holds.
All infinitesimal perturbations are annihilated in two steps, so no
Lyapunov exponent can be positive (with the usual extended-real
convention their growth rates are $-\infty$). This also directly
excludes the source's numerical notion of chaos.

Finally $|\alpha+1|>0=|\beta|$ because $\alpha\ne-1$.
For the explicit nonreal choice $\alpha=i$, the squared left-hand
side is $|1+i|^2=2$, and all admissible solutions have constant tail
$i,i,i,\ldots$. Thus the sufficiency direction, and hence the
claimed equivalence, is false.
\end{proof}

The counterexample concerns every admissible initial pair at the chosen
parameters, rather than merely a single fixed-point initial condition.
Undefined initial data at $z_0=-1$ do not generate solutions of the
original recurrence and cannot supply a chaotic orbit.

\section*{Four proposed follow-up tasks}
\begin{enumerate}
\item[Task 2.] Determine whether any sufficient chaos criterion depending
only on $|\alpha+1|$ and $|\beta|$ survives after imposing $\beta\ne0$.
\item[Task 3.] Find explicit nonzero parameter regions admitting a
forward-invariant domain on which all trajectories converge to one
fixed point, and compare them with the proposed inequality.
\item[Task 4.] Give rigorous invariant sets and certified Lyapunov
bounds for the numerical parameter pairs in the source's Table 4.
\item[Task 5.] Classify parameter pairs for which some iterate of the
state-space map is constant or has lower-dimensional image.
\end{enumerate}

\begin{thebibliography}{9}
\bibitem{Hassan} Sk.~Sarif Hassan,
\emph{Dynamics of $z_{n+1}=(\alpha+\alpha z_n+\beta z_{n-1})/(1+z_n)$
in Complex Plane},
\href{https://arxiv.org/abs/1511.04363v1}{arXiv:1511.04363v1} (2015),
p.~1 and Conjecture~5.1, p.~13.
\end{thebibliography}
\end{document}
